\documentclass{article}
\usepackage{amsmath, amssymb, amsthm}
\usepackage{hyperref}
\usepackage{geometry}
\geometry{margin=1in}

\title{Erd\H{o}s Problem \#1144}
\date{Last updated: 2026-01-23}
\author{Source: erdosproblems.com}

\begin{document}
\maketitle

\noindent\textbf{Status:} OPEN \quad \textbf{No prize}

\noindent\textbf{Tags:} number theory, probability

\medskip

\noindent\textbf{Problem Statement:}

Let $f$ be a random completely multiplicative function, where for each prime $p$ we independently choose $f(p)\in \{-1,1\}$ uniformly at random. Is it true that\[\limsup_{N\to \infty}\frac{\sum_{m\leq N}f(m)}{\sqrt{N}}=\infty\]with probability $1$?




This model of a random multiplicative function is sometimes called a Rademacher function, although this is sometimes reserved for a merely multiplicative function (which is $0$ on non-squarefree integers). See [520] for the partial sums of this alternative model. It should also be compared to another popular model of random completely multiplicative functions, Steinhaus functions, which have $f(p)$ uniformly distributed over the unit circle. Atherfold \cite{At25} has proved that, almost surely,\[\sum_{m\leq N}f(m)\ll N^{1/2}(\log N)^{1+o(1)}.\]




References

[At25] C. Atherfold, Almost sure bounds for weighted sums of Rademacher random multiplicative functions . arXiv:2501.11076 (2025).


\bigskip
\noindent\small{Source: \url{https://www.erdosproblems.com/1144}}
\end{document}
